% ============================================================
\section{Kiến trúc hệ thống}
\label{sec:arch}
% ============================================================

\subsection{Pipeline tổng thể}

\Cref{fig:pipeline} mô tả đường đi của một câu hỏi. Mỗi tầng là một module Python
riêng, nhận dữ liệu có kiểu tường minh từ tầng trước và ghi kết quả của mình vào một
file output.

\begin{figure}[H]
  \centering
  \begin{tikzpicture}[
      font=\small,
      stage/.style={draw, rounded corners=2pt, fill=blue!6, minimum width=5.4cm,
        minimum height=8.5mm, align=center},
      res/.style={draw, rounded corners=2pt, fill=orange!10, minimum width=3.6cm,
        minimum height=8.5mm, align=center, font=\footnotesize},
      io/.style={draw, rounded corners=6pt, fill=green!8, minimum width=5.4cm,
        minimum height=7mm, align=center},
      outfile/.style={font=\footnotesize\ttfamily, text=black!70, anchor=west},
      arr/.style={-{Stealth[length=2mm]}, thick},
      lab/.style={font=\scriptsize\itshape, text=black!65, right=1mm}
    ]
    \node[io] (in) {Câu hỏi tiếng Việt};
    \node[stage, below=6mm of in] (tok) {Tokenizer \\ \scriptsize chuẩn hóa NFC, chữ thường, tách dấu câu};
    \node[stage, below=6mm of tok] (parse) {Earley Parser \\ \scriptsize nhận/từ chối, dựng cây có nhãn};
    \node[stage, below=6mm of parse] (sem) {Semantic Interpreter \\ \scriptsize nhánh \texttt{Q\_*} $\rightarrow$ intent, slot};
    \node[stage, below=6mm of sem] (ctx) {Dialogue Context \\ \scriptsize giải tham chiếu; ghi nhớ thực thể khi FOUND};
    \node[stage, below=6mm of ctx] (query) {Query Engine \\ \scriptsize tra cứu theo intent/slot};
    \node[stage, below=6mm of query] (ans) {Answer Generator \\ \scriptsize template theo loại dữ kiện};
    \node[io, below=6mm of ans] (outn) {Câu trả lời + trace};

    \node[res, left=9mm of parse] (cfg) {\texttt{grammar.cfg} \\ CFG 26 nhánh};
    \node[res, left=9mm of sem] (lex) {Entity lexicon \\ \texttt{entities.txt}};
    \node[res, left=9mm of query] (kb) {Knowledge base \\ 6 file \texttt{data/kb/}};

    \foreach \a/\b/\t in {in/tok/{str}, tok/parse/{tokens}, parse/sem/{Tree hoặc ()},
        sem/ctx/{SemanticFrame}, ctx/query/{frame đã giải}, query/ans/{QueryResult},
        ans/outn/{}}
      {\draw[arr] (\a) -- (\b) node[midway, lab] {\t};}
    \draw[arr, dashed] (cfg) -- (parse);
    \draw[arr, dashed] (lex) -- (sem);
    \draw[arr, dashed] (kb) -- (query);
    \node[outfile] at ($(parse.east)+(1.2cm,0)$) {parse-results.txt};
    \node[outfile] at ($(sem.east)+(1.2cm,0.2)$) {semantic.txt};
    \node[outfile] at ($(sem.east)+(1.2cm,-0.2)$) {intent-entity.txt};
    \node[outfile] at ($(query.east)+(1.2cm,0)$) {query.txt};
    \node[outfile] at ($(ans.east)+(1.2cm,0)$) {answer.txt};
    \foreach \n in {parse, sem, query, ans}
      {\draw[black!35, densely dotted] (\n.east) -- ++(1.1cm,0);}
  \end{tikzpicture}
  \caption{Pipeline xử lý một câu hỏi và file output tương ứng}
  \label{fig:pipeline}
\end{figure}

Một lần chạy \code{run.py all} đi qua pipeline với toàn bộ 73 câu đầu vào và sinh đủ
9 file trong \code{output/}: ngoài 5 file theo từng câu ở \cref{fig:pipeline} còn có
\code{grammar.txt} (văn phạm dùng cho lần chạy), \code{samples.txt} (câu sinh từ văn
phạm), \code{dialogue.txt} (6 hội thoại) và \code{evaluation.txt} (kết quả đánh giá).

\subsection{Module và kiểu dữ liệu}

Mã nguồn nằm trong gói \code{hcmut.iaslab.nlp.app} (khoảng 3.000 dòng Python, chỉ dùng
thư viện chuẩn). \Cref{tab:modules} liệt kê các module chính.

\begin{table}[htbp]
  \centering
  \small
  \begin{tabularx}{\linewidth}{|P{3.2cm}|L|P{3.6cm}|}
    \hline
    \textbf{Module} & \textbf{Vai trò} & \textbf{Kiểu dữ liệu ra} \\ \hline
    \code{tokenizer.py} & Chuẩn hóa và tách token & \code{list[str]} \\ \hline
    \code{grammar.py} & Đọc \code{grammar.cfg}, kiểm tra cú pháp luật & \code{Grammar} \\ \hline
    \code{parser.py}, \code{tree.py} & Parser Earley, cây cú pháp có nhãn & \code{ParseResult}, \code{Tree} \\ \hline
    \code{generator.py} & Sinh câu hợp lệ từ văn phạm & \code{list[str]} \\ \hline
    \code{semantic.py}, \code{entities.py} & Diễn giải cây, chuẩn hóa entity về ID & \code{SemanticFrame} \\ \hline
    \code{dialogue.py} & Bộ nhớ và giải tham chiếu cục bộ & \code{SemanticFrame} đã giải \\ \hline
    \code{kb.py} & Nạp và kiểm tra KB & \code{KnowledgeBase} \\ \hline
    \code{query.py} & Truy vấn KB theo intent/slot & \code{QueryResult} \\ \hline
    \code{answer.py} & Sinh câu trả lời và fallback & \code{str} \\ \hline
    \code{pipeline.py} & Ghép các tầng cho một câu & \code{PipelineResult} \\ \hline
    \code{evaluate.py} & Chấm các bộ kiểm thử & báo cáo đánh giá \\ \hline
    \code{cli.py}, \code{output.py} & Dòng lệnh, ghi 9 file output & file văn bản \\ \hline
  \end{tabularx}
  \caption{Module chính của hệ thống}
  \label{tab:modules}
\end{table}

Ba kiểu dữ liệu nối các tầng:
\begin{itemize}
  \item \code{SemanticFrame}: \code{intent}, \code{entity}, \code{slots} (từ điển
        slot--giá trị) và \code{source_branch} (nhánh \texttt{Q\_*} sinh ra frame).
  \item \code{QueryResult}: \code{found}, \code{kind} (loại dữ kiện, ví dụ
        \code{SCHEDULE_WEEK}), \code{data}, \code{sources} (file KB đã dùng) và
        \code{reason} khi không tìm thấy.
  \item \code{PipelineResult}: câu gốc, token, cây, frame gốc, frame sau khi giải
        ngữ cảnh, truy vấn, nguồn và câu trả lời; đây là nội dung của chế độ
        \code{--trace}.
\end{itemize}

\subsection{Ranh giới giữa các tầng}
\label{sec:contracts}

Nhóm viết hợp đồng giữa các tầng trước khi chia việc hiện thực. Mỗi hợp đồng quy
định đầu vào, đầu ra và bất biến của một tầng; các ràng buộc quan trọng:
\begin{itemize}
  \item Tầng ngữ nghĩa \textbf{không đọc KB}: intent và entity chỉ được suy ra từ cây
        cú pháp và entity lexicon.
  \item Tầng truy vấn không đọc lại câu gốc, chỉ dùng frame ngữ nghĩa.
  \item Tầng trả lời chỉ dùng \code{QueryResult}; một kết quả \code{found=True} phải
        sinh câu trả lời chứa dữ kiện, còn \code{found=False} phải đi vào fallback.
\end{itemize}
Mỗi hợp đồng có một bộ test nghiệm thu dùng đầu vào golden, nên từng tầng kiểm thử
được độc lập với tầng trước. Bộ test hiện có 49 test, gồm test hợp đồng cho từng tầng
và các ca hồi quy (\cref{sec:testing}).
